% =============================================================================
% Appendix A: Architecture Variations
% =============================================================================
\chapter{Architecture Variations}
\label{app:architecture-variations}

This appendix surveys architectural choices not used in \TryLM but important to know.

% -----------------------------------------------------------------------------
\section{Multi-Query Attention (MQA)}
% -----------------------------------------------------------------------------

\textbf{Standard multi-head}: Each head has its own Q, K, V projections.

\textbf{MQA}: All heads share K and V; only Q is per-head.

\begin{table}[H]
\centering
\begin{tabular}{lll}
\toprule
\textbf{Approach} & \textbf{KV Cache Size} & \textbf{Quality} \\
\midrule
Multi-head (TryLM) & $h \times d_k$ & Best \\
Multi-query (MQA) & $d_k$ & Slightly lower \\
Grouped-query (GQA) & $g \times d_k$ & Middle ground \\
\bottomrule
\end{tabular}
\caption{Attention variants comparison}
\end{table}

\textbf{Use case}: MQA/GQA reduce memory during inference, important for large-scale deployment.

% -----------------------------------------------------------------------------
\section{Flash Attention}
% -----------------------------------------------------------------------------

Standard attention requires $O(n^2)$ memory for the attention matrix. Flash Attention:
\begin{itemize}
    \item Computes attention in tiles
    \item Never materializes full attention matrix
    \item Achieves $O(n)$ memory
    \item 2--4x faster for long sequences
\end{itemize}

\begin{lstlisting}[style=python]
# PyTorch 2.0+ includes Flash Attention
from torch.nn.functional import scaled_dot_product_attention

# Automatically uses Flash Attention when available
output = scaled_dot_product_attention(q, k, v, is_causal=True)
\end{lstlisting}

% -----------------------------------------------------------------------------
\section{Mixture of Experts (MoE)}
% -----------------------------------------------------------------------------

MoE routes each token to a subset of ``expert'' FFN layers:

\begin{itemize}
    \item \textbf{More parameters}: 8--64 experts, only 1--2 active per token
    \item \textbf{Similar compute}: Only active experts run
    \item \textbf{Scaling}: Enables very large models (Mixtral, GPT-4)
\end{itemize}

\begin{lstlisting}[style=python]
class MoELayer(nn.Module):
    def __init__(self, n_experts=8, top_k=2):
        self.experts = nn.ModuleList([FFN() for _ in range(n_experts)])
        self.router = nn.Linear(d_model, n_experts)

    def forward(self, x):
        # Route each token to top-k experts
        routing_weights = F.softmax(self.router(x), dim=-1)
        top_k_weights, top_k_indices = routing_weights.topk(2)
        # Compute weighted sum of expert outputs
        ...
\end{lstlisting}

% -----------------------------------------------------------------------------
\section{Alternative Position Encodings}
% -----------------------------------------------------------------------------

\textbf{ALiBi} (Attention with Linear Biases):
\begin{itemize}
    \item Adds linear bias based on distance: $\text{bias}_{ij} = -m \cdot |i - j|$
    \item No learned parameters
    \item Excellent length extrapolation
\end{itemize}

\textbf{Relative positional encodings}:
\begin{itemize}
    \item T5-style: Learned biases based on relative distance
    \item Bucketed distances for efficiency
\end{itemize}

% -----------------------------------------------------------------------------
\section{Encoder-Decoder Architectures}
% -----------------------------------------------------------------------------

\TryLM is decoder-only. Alternative architectures:

\begin{description}
    \item[BERT (Encoder-only)] Bidirectional attention, used for understanding tasks
    \item[T5 (Encoder-Decoder)] Separate encoder and decoder, used for sequence-to-sequence
    \item[GPT (Decoder-only)] Causal attention, used for generation (TryLM)
\end{description}

\begin{figure}[H]
\centering
\begin{tikzpicture}[
    box/.style={rectangle, draw, rounded corners, minimum width=2cm, minimum height=1cm, fill=blue!10},
]
    % BERT
    \node[box] (bert) at (0, 0) {Encoder};
    \node[below=0.2cm of bert, font=\small] {BERT};
    \node[above=0.2cm of bert, font=\footnotesize] {Bidirectional};

    % T5
    \node[box] (t5e) at (4, 0.5) {Encoder};
    \node[box] (t5d) at (4, -0.5) {Decoder};
    \draw[-{Stealth}] (t5e) -- (t5d);
    \node[below=0.2cm of t5d, font=\small] {T5};

    % GPT
    \node[box] (gpt) at (8, 0) {Decoder};
    \node[below=0.2cm of gpt, font=\small] {GPT/TryLM};
    \node[above=0.2cm of gpt, font=\footnotesize] {Causal};
\end{tikzpicture}
\caption{Transformer architecture variants}
\end{figure}
